\section{本讲小结}
开关题由稳态确定储能初态，再按新通路求暂态。电感电流连续、电容电压连续均需相应无冲激条件；灯泡亮度由自身瞬时功率决定。

一阶电路使用新连接下的时间常量；总电量先判断积分是否收敛。LC与电磁机械振荡都须由方程及两个初态求相位，含电阻时再检查阻尼与能量下降。


快切换先判断是否存在冲激，再连接快过程前后的状态，有限热量不能因持续时间短而丢弃。含多个储能元件时，应由实际连接确定独立组合；电阻可能只耗散部分组合，未经过电阻的环流仍可保留磁能。
